\section{Construction on a moving ball}\label{sec:construction}

We fix
\[
 a=\frac12,\qquad A=\frac1{16},\qquad U=100,\qquad
 S=1,\qquad R_0=\frac{101}{100}.
\]
The small positive order $\alpha$ will be chosen below. We write
$s=1-\alpha/2$, $\beta=1-\alpha$, and
\[
 L_\alpha=(-\Delta)^{\alpha/2},\qquad
 K_\alpha=(-\Delta)^{-1+\alpha/2},\qquad
 v[w]=-\nabla K_\alpha w,\qquad f(k)=k^a.
\]
Thus the equation is $u_t+\nabla\cdot(v[u]f(u))=0$ and
$\nabla\cdot v[u]=L_\alpha u$. Densities defined on a ball are always
extended by zero when a whole-space operator is applied.

\subsection{The profile and the pressure field}

Set $h(z)=10z^3-15z^4+6z^5$ on $[0,1]$, and define the radial datum
\begin{equation}\label{eq:datum}
 F(r)=\begin{cases}
 1,&0\leq r<S,\\
 U-(U-A)h\bigl((r-S)/(R_0-S)\bigr),&S<r<R_0,\\
 0,&r>R_0.
 \end{cases}
\end{equation}
Its values on the two spheres are immaterial. There is an upward jump from
$1$ to $U$ at $S$, and a downward jump from $A$ to zero at $R_0$.
The first two derivatives of the annular profile vanish at both endpoints.
In particular $F$ is nonzero, nonnegative, compactly supported and in
$BV\cap L^1\cap L^\infty$, with $\|F\|_\infty=U$.
Every positive power of $F$ belongs to $H^{\alpha/2}(\R^3)$ for $0<\alpha<1$.
Indeed it is piecewise Lipschitz, and flattening a jump sphere bounds the
cross-interface part of its Gagliardo integral by a constant times
\[
 \int_0^1\int_0^1(d_1+d_2)^{-1-\alpha}\,\dd d_1\dd d_2<\infty.
\]
The same-region part follows from Lipschitz continuity, and the part away
from the diagonal follows from bounded support.

For $0\leq\alpha<1$ the normalized potential and velocity are
\begin{align}
 K_\alpha w(x)&=d_\alpha\int_{\R^3}|x-y|^{-1-\alpha}w(y)\,\dd y,
 &d_\alpha&=\frac{\Gamma((1+\alpha)/2)}
 {2^{2-\alpha}\pi^{3/2}\Gamma(1-\alpha/2)},\label{eq:riesz}\\
 v[w](x)&=c_\alpha\int_{\R^3}
 \frac{x-y}{|x-y|^{3+\alpha}}w(y)\,\dd y,
 &c_\alpha&=(1+\alpha)d_\alpha.\label{eq:velocity}
\end{align}
These integrals are continuous for bounded integrable densities.
The velocity kernel is locally integrable since
$\int_0^1r^{-\alpha}\,\dd r<\infty$. Its near-field contribution on a ball
of radius $\delta$ is at most $K\|w\|_\infty\delta^{1-\alpha}$;
the corresponding potential contribution is at most
$K\|w\|_\infty\delta^{2-\alpha}$. The far kernels are bounded, so mass
controls the remaining contributions. For fixed $0<\alpha<1$, splitting a
velocity difference at twice the displacement gives
\begin{equation}\label{eq:velocity-holder}
 |v[w](x)-v[w](x')|\leq K\|w\|_\infty |x-x'|^{1-\alpha}
\end{equation}
on a common bounded support, with the far-field contribution controlled by
mass. Constants in this estimate need not stay bounded as $\alpha\downarrow0$.

For clarity, the Fourier normalization in \eqref{eq:riesz} can be checked
without an operator convention left implicit. Insert
\[
 |z|^{-1-\alpha}=\frac1{\Gamma((1+\alpha)/2)}
 \int_0^\infty t^{(1+\alpha)/2-1}e^{-t|z|^2}\,\dd t.
\]
The unitary Fourier transform of the Gaussian is
$(2t)^{-3/2}e^{-|\xi|^2/(4t)}$. Substitution of $|\xi|^2/(4t)$,
including the convolution factor $(2\pi)^{3/2}$, gives the multiplier
$|\xi|^{\alpha-2}$. Testing first on Schwartz functions and then truncating
the locally integrable kernels justifies the identity in distributions.
Differentiation gives \eqref{eq:velocity} and
\begin{equation}\label{eq:divv}
 -\Delta K_\alpha w=\nabla\cdot v[w]=L_\alpha w.
\end{equation}
The singular-integral normalization, in dimension $n=1$ or $3$, is
\begin{equation}\label{eq:frac-kernel}
 L_{n,\alpha}w(x)=c_{n,\alpha}\,\operatorname{PV}
 \int\frac{w(x)-w(y)}{|x-y|^{n+\alpha}}\,\dd y,
 \qquad c_{n,\alpha}=
 \frac{\alpha2^{\alpha-1}\Gamma((n+\alpha)/2)}
 {\pi^{n/2}\Gamma(1-\alpha/2)}.
\end{equation}
One obtains this coefficient by the same Gaussian representation applied to
$1-\cos(z\cdot\xi)$ and one integration by parts. Symmetric truncation and
Plancherel identify the positive quadratic form with
$\|(-\Delta)^{\alpha/4}w\|_2^2$.

Two further bounds will be used repeatedly. Inside a ball $D$,
\begin{equation}\label{eq:indicator}
 0\leq L_\alpha\mathbf1_D(x)\leq K\operatorname{dist}(x,\partial D)^{-\alpha}.
\end{equation}
Also $L_\alpha$ maps bounded compactly supported $C^\theta$ functions to
bounded functions if $\theta>\alpha$. Both statements follow by splitting
the singular integral at the distance to the boundary or at one. The
radial integrals are respectively $\int_d^\infty r^{-1-\alpha}\dd r$ and
$\int_0^1r^{\theta-\alpha-1}\dd r$.

For radial $w$ write $v[w](r)$ for the outward component, and put
\[
 W[w;r]=4\pi r^2v[w](r).
\]
If $w\geq0$ is supported in $B_r$ and has mass $M>0$, then at $x=rn$,
$|n|=1$, the inequality $r-y\cdot n\geq |rn-y|^2/(2r)$ gives
\begin{equation}\label{eq:boundary-positive}
 v[w](r)\geq c_\alpha M(2r)^{-2-\alpha}>0.
\end{equation}
At $\alpha=0$, Gauss' formula yields
$v_0[w](r)=r^{-2}\int_0^r\rho^2w(\rho)\,\dd\rho$.
Consequently
\[
 W_0[F;S]=\frac{4\pi}3,\qquad W_0[F;R_0]=\int F,
 \qquad
 \frac{W_0[F;R_0]}{W_0[F;S]}
 \leq1+100\bigl((101/100)^3-1\bigr)=4.0301.
\]
For this fixed datum the two velocity values converge to their Newtonian
values as $\alpha\downarrow0$. A majorant with any fixed order less than
one controls the near kernel, and dominated convergence applies on the
remaining compact support. Thus both weights are positive and their ratio
is less than five for all sufficiently small positive $\alpha$.

The concavity deficits that will be compared are
\begin{align}
 g_A(k)&=k^a-A^{a-1}k, &0\leq k\leq A,\\
 H_U(k)&=k^a-1-\frac{U^a-1}{U-1}(k-1), &1\leq k\leq U.
\end{align}
They are positive in the interiors by strict concavity. Direct integration gives
\begin{equation}\label{eq:margin}
 G_A:=\int_0^A\frac{g_A(k)}k\,\dd k=\frac14,
 \qquad J:=\int_1^U\frac{H_U(k)}k\,\dd k
 =9-\frac{10}{11}\log100>4.
\end{equation}
We henceforth fix one $0<\alpha<1/4$ sufficiently small that
\begin{equation}\label{eq:strict-margin}
 W[F;R_0]G_A<W[F;S]J.
\end{equation}
There is no numerical endpoint asserted for this choice. Notice also that,
for every $q>1$,
\begin{equation}\label{eq:density-cost}
 \int_0^A k^{q-2}g_A(k)\,\dd k\leq A^{q-1}G_A.
\end{equation}

\subsection{A viscous free boundary problem}

We approximate only the internal jump of $F$. There are radial
$g_j\in C^2(\overline B_{R_0})$, $A\leq g_j\leq U$, with
$g_j=A$, $\partial_ng_j=0$ on the boundary, whose zero extensions converge
to $F$ in every finite $L^q$ and in $L^1$. They agree with $F$ in one fixed
outer collar, with a common $C^1$ bound there. To construct them, mollify
the radial profile near $S$ and join it to $F$ with a fixed smooth cutoff.
Only values in $[A,U]$ are averaged. Their masses lie in a common interval
$[M_-,M_+]$ with $M_->0$, after deleting a finite prefix if necessary.
Choose $\eps_j=2^{-j}$, with a harmless relabeling after that deletion.

For fixed $\eps>0$ and one such datum consider, in $D(t)=B_{R(t)}$,
\begin{align}
 w_t&=\nabla\cdot(w^a\nabla K_\alpha w)+\eps\Delta w,
 &w|_{\partial D(t)}&=A,\label{eq:viscous}\\
 R'&=A^{a-1}v[w](R)-\frac\eps A\partial_nw,
 &R(0)&=R_0,\qquad w(0)=g.\label{eq:stefan}
\end{align}
The positive boundary value permits a discontinuity to exterior vacuum.
The law \eqref{eq:stefan} is precisely the cancellation required by mass
conservation for the zero extension.

We first give the local realization used in subsequent integrations by
parts. Choose $3<p_0<1/\alpha$. On $L^{p_0}(B_1)$ let
\[
 Gz=\Delta z,\qquad D(G)=W^{2,p_0}(B_1)\cap W^{1,p_0}_0(B_1),
 \qquad \Sigma(t)=e^{tG}.
\]
The Dirichlet Laplacian is densely defined and sectorial
\cite[Section~2.4]{Lunardi2004}. Its analytic semigroup satisfies
\begin{align}
 \|\Sigma(t)h\|_{C^1}&\leq K(1+t^{-\lambda})\|h\|_{p_0},
 &\lambda&=\frac12+\frac3{2p_0}<1,\label{eq:heat-smoothing}\\
 \|[\Sigma(h)-I]\Sigma(r)\|_{L^{p_0}\to C^1}
 &\leq K\min(r^{-\lambda},hr^{-1-\lambda}).\label{eq:heat-increment}
\end{align}
Here and below bounded time intervals are understood in heat estimates.
The semigroup is uniformly bounded and strongly continuous on the space of
$C^1$ functions with zero boundary value, in the full $C^1$ norm.
This space imposes no zero normal derivative.

Here are details of the boundary and $C^1$ assertions. Comparison with the
whole-space Gaussian gives the $L^{p_0}$ to $L^\infty$ estimate. At a tangent
plane to the convex ball, comparison with
$\|h\|_\infty\operatorname{erf}(\ell/(2\sqrt t))$, where $\ell$ is the
inward affine distance, bounds the boundary gradient by
$\|h\|_\infty/\sqrt{\pi t}$. For a heat solution,
\[
 (\partial_t-\Delta)(t|\nabla z|^2+z^2)
 =-|\nabla z|^2-2t|D^2z|^2.
\]
The maximum principle and semigroup factorization give
\eqref{eq:heat-smoothing}. For zero-trace $C^1$ data with Lipschitz constant
$L$, the stationary barriers $\pm L\ell$ instead give a uniform gradient
bound. For smooth such data,
$\Sigma(t)h-h=\int_0^t\Sigma(r)\Delta h\,\dd r$ tends to zero in $C^1$
by $\lambda<1$. Smooth zero-trace approximation and the uniform bound
extend this to all zero-trace $C^1$ data. Finally the analytic estimate
$\|G\Sigma(r)\|_{L^{p_0}\to L^{p_0}}\leq K/r$
\cite[Theorem~1.2.3]{Lunardi2004}, combined with
\eqref{eq:heat-smoothing}, proves \eqref{eq:heat-increment} by integration.

To solve \eqref{eq:viscous}--\eqref{eq:stefan}, put
$y=x/R(t)$, $\dd\tau/\dd t=\eps/R(t)^2$, $\rho(\tau)=R(t)$ and
$z(\tau,y)=w(t,x)$. Write
$\theta=\rho^{2-\alpha}/\eps$ and
$\widetilde v=v[z\mathbf1_{B_1}]$. The equations become
\begin{align}
 z- A&=\Sigma(\tau)(g(R_0\cdot)-A)
       +\int_0^\tau\Sigma(\tau-r)\mathcal N(z(r),\rho(r))\,\dd r,
                                                        \label{eq:mild}\\
 \rho_\tau&=\theta\rho A^{a-1}\widetilde v(1)
              -\frac\rho A\partial_nz(1),\label{eq:rho}\\
 \mathcal N(z,\rho)&=
 \left[\theta A^{a-1}\widetilde v(1)y
 -\frac{\partial_nz(1)}A y-a\theta z^{a-1}\widetilde v(y)\right]\cdot\nabla z
 \nonumber\\[-2mm]
 &\hspace{8mm}-\theta z^a
 \left(L_\alpha[(z-A)\mathbf1_{B_1}]+A L_\alpha\mathbf1_{B_1}\right).
                                                        \label{eq:forcing}
\end{align}
For positive radial $C^1$ functions with trace $A$ and positive radii,
$\mathcal N$ is locally Lipschitz into $L^{p_0}$. Indeed the zero extension
of $z-A$ is globally Lipschitz, so its fractional Laplacian is bounded;
\eqref{eq:indicator} is in $L^{p_0}$ because $\alpha p_0<1$.
All boundary derivatives and the remaining products are locally Lipschitz
in the indicated $C^1$ norm. The radius vector field has the same property.
On continuous $C^1$ density paths and continuous radius paths, the two
integral equations have contraction constants bounded by
$K\tau_0^{1-\lambda}$ and $K\tau_0$. Strong $C^1$ continuity of the linear
path permits a positive closed ball of paths on a sufficiently short
interval. The contraction gives a unique positive radial local solution.

This mild solution has the required strict realization. The increment
estimate \eqref{eq:heat-increment}, split at $r=h$, gives a time
$C^{1-\lambda}$ bound in $C^1$ for the Duhamel term. The initial datum
$g(R_0\cdot)-A$ belongs to $D(G)$, and the same bound follows for its heat
path by integrating $\Sigma(r)G(g(R_0\cdot)-A)$. Hence
$\mathcal N(z,\rho)$ is H\"older continuous in time into $L^{p_0}$.
The strict-solution theorem \cite[Theorem~4.2.6(ii)]{Lunardi2004} applies:
the compatibility condition is
$G(g(R_0\cdot)-A)+\mathcal N(z(0),R_0)\in\overline{D(G)}$,
which is automatic since $D(G)$ is dense in $L^{p_0}$. Thus
\[
 z-A\in C^1_\tau L^{p_0}\cap C_\tau D(G).
\]
Interior localization and the Gaussian heat formula give classical
$C^1$ time and $C^2$ space regularity for positive times: after subtracting
the forcing value at the evaluation point, the differentiated heat kernel
has an integrable time singularity by the positive H\"older exponent.
The established $W^{2,p_0}$ realization gives spatial
$C^{1,1-3/p_0}$ regularity by Sobolev embedding. Thus the drift--gradient
product is H\"older on interior subballs, and the fractional kernel
estimates also make the remaining forcing terms spatially H\"older.
This also treats the origin in Cartesian
coordinates. Finally
\[
 v[w](Ry)=R^{1-\alpha}\widetilde v(y),\qquad
 L_\alpha w(Ry)=R^{-\alpha}L_\alpha[z\mathbf1_{B_1}](y),
 \qquad \frac{\dd t}{\dd\tau}=\frac{\rho^2}\eps
\]
undo the transformation and give exactly
\eqref{eq:viscous}--\eqref{eq:stefan}. Regularity here concerns the interior
and its one-sided boundary traces; the exterior jump is retained.

\subsection{Identities at the moving boundary}

Put $j=w^a\nabla K_\alpha w+\eps\nabla w$ inside $D(t)$.
The boundary law says $AR'=-j\cdot n$. Reynolds' rule therefore gives
\begin{equation}\label{eq:extended-mass}
 \partial_tw=\nabla\cdot(j\mathbf1_{D(t)}),\qquad
 \int w(t)=\int g
\end{equation}
for the zero extension, including its initial trace. At an interior global
maximum, $L_\alpha w\geq0$, $\nabla w=0$ and $\Delta w\leq0$.
The equation and a strict first-contact argument imply $w\leq U$ as long
as the positive local solution exists.

For an entropy $\eta$ with $\eta(0)=0$, define
\begin{equation}\label{eq:entropy-transforms}
 \Psi_\eta(k)=\int_0^k\eta''(z)f(z)\,\dd z,
 \quad F_\eta(k)=\int_0^k\Psi_\eta(z)\,\dd z,
 \quad Q_\eta(k)=\eta'(k)f(k)-\Psi_\eta(k).
\end{equation}
For local distributional identities we normalize also $\eta'(0)=0$.
The powers $\eta(k)=k^q$, $q>1$, and $k\log k$ are admitted by
approximation whenever the displayed integrals at zero converge.
For increasing $G$ set
\begin{equation}\label{eq:pair}
 \mathcal B_\alpha(G,w)=\frac{c_{3,\alpha}}2
 \iint\frac{(G(w(x))-G(w(y)))(w(x)-w(y))}
 {|x-y|^{3+\alpha}}\,\dd x\dd y.
\end{equation}
This is a nonnegative extended integral. When finite it is the symmetric
form of $\int G(w)L_\alpha w$.

Reynolds' rule and integration by parts yield the exact identity
\begin{align}
 \frac\dd{\dd t}\int_D\eta(w)
 +\mathcal B_\alpha(\Psi_\eta,w)
 +\eps\int_D\eta''(w)|\nabla w|^2
 ={}&W[w;R]\int_0^A\eta''(k)g_A(k)\,\dd k\nonumber\\
 &+\eps|\partial B_R|\partial_nw
       \left(\eta'(A)-\frac{\eta(A)}A\right).
                                                        \label{eq:entropy-global}
\end{align}
Here $W[w;R]=4\pi R^2v[w](R)$ and
$g_A(k)=k^a-A^{a-1}k$.
In fact the initial boundary term is
$j\cdot n[\eta'(A)-\eta(A)/A]$. Since
$\nabla\Psi_\eta(w)=\eta''(w)w^a\nabla w$ and
$\nabla\cdot\nabla K_\alpha w=-L_\alpha w$, integrating the pressure
term contributes $W[w;R]\Psi_\eta(A)$.
The remaining coefficient is
\[
 \Psi_\eta(A)-A^a\eta'(A)+A^{a-1}\eta(A)
 =\int_0^A\eta''(k)g_A(k)\,\dd k.
\]
This proves \eqref{eq:entropy-global}, including its sign. For fixed
$\eps$ the established strict $L^{p_0}$ realization, positive density and
$C^1$ traces justify the time-integrated calculation; the boundary
singularity has size $d^{-\alpha}$ and is integrable. Entropy approximation
on positive interior values gives the asserted non-smooth-at-zero cases.
For powers,
\begin{equation}\label{eq:psi-power}
 \Psi_{k^q}(k)=\frac{q(q-1)}{a+q-1}k^{a+q-1}.
\end{equation}
For the logarithmic entropy, \eqref{eq:entropy-global} reads
\begin{equation}\label{eq:log-identity}
 \frac\dd{\dd t}\int w\log w+
 \frac1a\int w^aL_\alpha w+
 \eps\int_D\frac{|\nabla w|^2}{w}
 =W[w;R]G_A+\eps|\partial B_R|\partial_nw.
\end{equation}

The pressure energy is
\[
 \mathcal E_\alpha(w)=\frac12\|(-\Delta)^{-s/2}w\|_2^2
 =\frac12\int wK_\alpha w.
\]
It is finite for bounded compact densities, since the potential kernel is
locally integrable and bounded away from its singularity. Its identity is
\begin{equation}\label{eq:pressure-identity}
 \frac\dd{\dd t}\mathcal E_\alpha(w)
 +\int_Dw^a|\nabla K_\alpha w|^2
 +\eps\int wL_\alpha w=\eps A W[w;R].
\end{equation}
To justify differentiating the energy, truncate the symmetric double
potential integral and differentiate both moving domains. The terms with
$AR'$ cancel those with $j\cdot n$, leaving
$-\int_D\nabla K_\alpha w\cdot j$.
The potential and its first derivative have integrable majorants;
the potential is also integrable on the two-dimensional boundary since
$\alpha<1$. Truncations can thus be removed after time integration.
Finally
$\int_D\nabla w\cdot\nabla K_\alpha w
=\int wL_\alpha w-AW[w;R]$, proving
\eqref{eq:pressure-identity}.

\subsection{Uniform barriers and continuation}

Recall $\beta=1-\alpha$. The regularized boundary layer must respond to the singularity
$d^{-\alpha}$ in \eqref{eq:indicator}. For $\kappa>0$ define
\begin{equation}\label{eq:layer}
 \ell_\eps(d)=\int_0^d y_\eps(z)\,\dd z,
 \qquad y_\eps(d)=\frac1\eps\int_0^\infty
 e^{-\kappa h/\eps}(d+h)^{-\alpha}\,\dd h.
\end{equation}
The function is increasing and concave, and
\begin{align}
 \kappa\ell_\eps'-\eps\ell_\eps''&=d^{-\alpha},
 &\ell_\eps(d)&\leq\frac{d^\beta}{\kappa\beta},\label{eq:layer-eq}\\
 \ell_\eps'(0)&=\Gamma(1-\alpha)\kappa^{\alpha-1}\eps^{-\alpha}.
                                                        \label{eq:layer-slope}
\end{align}
These follow by differentiation and the Gamma integral. The decreasing
derivative gives
$|\ell_\eps(d+h)-\ell_\eps(d)|\leq\ell_\eps(|h|)$,
so the $C^\beta$ seminorm is uniform. The exponential approximate identity
in \eqref{eq:layer} gives uniform convergence on bounded intervals to
$d^\beta/(\kappa\beta)$, by the integrable bound in
\eqref{eq:layer-eq} and this common modulus.

\begin{lemma}\label{lem:uniform-family}
After deleting a finite prefix of the seed sequence, there is a common
$T>0$ on which \eqref{eq:viscous}--\eqref{eq:stefan} have the local
realizations above. Constants independent of $j$ give
\begin{gather}
 A/2\leq w_j\leq U\quad\hbox{in }D_j(t),\qquad
 R_0/2<R_j(t)<2R_0,\qquad |R_j'|\leq K,\label{eq:uniform-basic}\\
 \max\{c(t),A-K\ell_{\eps_j}(R_j-r)\}
 \leq w_j(t,r)\leq
 \min\{U,A+K\ell_{\eps_j}(R_j-r)\},\label{eq:barriers}\\
 |\eps_j\partial_nw_j|\leq K\eps_j^\beta.\label{eq:boundary-flux}
\end{gather}
Here $c(0)=3A/4$ and $A/2\leq c(t)\leq3A/4$.
Constants may depend on the fixed $\alpha$, profile and seed collar bounds,
but not on the interior seed gradients.
\end{lemma}

\begin{proof}
Initially stop a regular solution when its radius leaves $[R_0/2,2R_0]$
or $|\eps\partial_nw|$ reaches a small constant $J_0$.
Mass conservation, the cap, and \eqref{eq:boundary-positive} give
$v[w](R)\geq v_*>0$, uniformly before this stopping time.
The velocity is uniformly bounded and has the common spatial modulus
\eqref{eq:velocity-holder}. For a collar state $z\geq A/2$,
\begin{align*}
 az^{a-1}v[w](r)-R'
 \leq{}&-[1-a2^{1-a}]A^{a-1}v[w](R)\\
 &+a(A/2)^{a-1}|v[w](r)-v[w](R)|+J_0/A.
\end{align*}
The margin $1-a2^{1-a}=1-1/\sqrt2$ is positive. Choose the collar width
$\delta_0$ and then $J_0$ so that the last two errors each take at most
one quarter of this margin. For some $\kappa>0$ this gives
\begin{equation}\label{eq:collar-drift}
 az^{a-1}v[w](r)-R'\leq-2\kappa
 \quad (R-\delta_0<r<R).
\end{equation}

Consider the two functions in \eqref{eq:barriers}. The lower one is initially
below $g_j\geq A$. A sufficiently large $K$ makes the upper one dominate
the common outer $C^1$ collar, and outside that collar it is the cap $U$.
This uses the uniform positive lower bound for $\ell_\eps$ at each fixed
positive distance and its concavity, which bound $d/\ell_\eps(d)$ there.
Subtract $A\mathbf1_D$ from either zero-extended barrier. The remainder
is globally $C^\beta$ with norm at most a constant times $K$.
Since $\beta>\alpha$, \eqref{eq:indicator} and the H\"older estimate for
$L_\alpha$ give
\[
 |L_\alpha(\hbox{barrier}-A\mathbf1_D)|\leq K_2K,
 \qquad 0\leq L_\alpha\mathbf1_D\leq K_2d^{-\alpha}.
\]
At a smooth contact in the lower collar, \eqref{eq:collar-drift} and
\eqref{eq:layer-eq} bound the residual of
$\partial_t+a z^{a-1}v[w]\cdot\nabla+z^aL_\alpha-\eps\Delta$
by
\[
 -K d^{-\alpha}+K_2U^aA d^{-\alpha}+K_2U^aK.
\]
The radial curvature term is favorable. First enlarge $K$ to absorb the
$A$ term, then reduce $\delta_0$ so that $d^{-\alpha}$ absorbs the last
term, and finally enlarge $K$ to place both nonconstant collars inside the
reduced $\delta_0$. The fractional bound is linear in $K$, so the last
choice preserves the sign. The upper collar has reversed favorable
transport and layer signs and is a supersolution by the same choices.

The constant lower branch lies a fixed positive distance from the boundary:
at its join, $A-c\geq A/4$ and
$\ell_\eps(d)\leq d^\beta/(\kappa\beta)$. At a minimum of this barrier
only exterior vacuum makes a positive contribution to its fractional
Laplacian; it is at most $K_3c$ with a uniform $K_3$.
Choose $c'=-K_1c$ to dominate this term. Then $c\geq A/2$ on a common
positive time interval. The upper constant branch is a global maximum.
The lower join is convex and the upper join concave, so the incompatible
one-sided slopes preclude a forbidden first contact there.

Comparison uses the actual field $v[w]$ throughout. At a lower smooth
contact the zero-extended difference has global minimum zero and
$L_\alpha(w-w_-)\leq0$; at an upper contact the sign is reversed.
For fixed $\eps$, the difference equation has a zeroth-order coefficient
bounded below: the barrier gradient is bounded at that $\eps$, its density
is positive, and its fractional Laplacian is bounded below by $-K_2K$.
An exponential time factor therefore gives the non-strict minimum
principle. This proves comparison without assuming monotonicity of the
map $w\mapsto v[w]$.

Dividing the collar inequalities by $d$ at the boundary and using
\eqref{eq:layer-slope} gives \eqref{eq:boundary-flux}, which is strictly
below $J_0$ for sufficiently small viscosities. Thus the flux stopping
condition cannot occur first. Equation \eqref{eq:stefan} bounds $R'$;
shortening the common interval prevents either radius exit or $c=A/2$.

It remains to exclude an earlier loss of the fixed-$\eps$ realization.
Recall the fixed-ball variables $y=x/R(t)$,
$z(\tau,y)=w(t,R(t)y)$, $\rho(\tau)=R(t)$ and
$\dd\tau/\dd t=\eps/R(t)^2$.
The boundary derivative in \eqref{eq:forcing} is bounded at fixed $\eps$
by \eqref{eq:boundary-flux}. The other drift coefficients are bounded by
the density cap, positive floor and radius bounds. The fractional map of
$(z-A)\mathbf1_{B_1}$ is linear and bounded from $C^1$ to $L^\infty$.
Consequently
\[
 \|\mathcal N(z,\rho)\|_{p_0}\leq K_\eps(1+\|z\|_{C^1}).
\]
Equation \eqref{eq:heat-smoothing} gives a Volterra inequality with kernel
$1+r^{-\lambda}$, $\lambda<1$. Split each bounded rescaled interval into
finitely many pieces on which
$K_\eps\int_0^\delta(1+r^{-\lambda})\dd r<1/2$;
successive absorption bounds the $C^1$ norm on the whole interval.
The contraction can be restarted from positive $C^1$ data using these same
bounds and heat comparison. At positive elapsed times, the increment
estimate gives H\"older forcing. Applying $G$ to Duhamel after subtracting
the terminal forcing value is legitimate since
$\int_0^h r^{-1}r^\theta\dd r<\infty$ for $\theta>0$; the constant part
has generator $(\Sigma(h)-I)\mathcal N(z(t),\rho(t))$.
Closedness of $G$ gives the strict realization there, and interior heat
localization gives its classical derivatives. Local uniqueness matches
this restart to the earlier strict solution. Finite physical time is
bounded rescaled time at fixed $\eps$ because the radius stays between
positive bounds. Hence each fixed-$\eps$ solution continues to the common
physical time. No $C^1$ bound uniform in $\eps$ is asserted.
\end{proof}
